\documentclass[../main.tex]{subfiles}

\begin{document}

\chapter{C++ Basics for Competitive Programming}

Before learning algorithms and solving problems, we need to choose
a language to use. There are many programming languages that we can
choose from for competitive programming such as C++, Python, and Java.
However, many competitors use C++ due to its speed and popularity
\cite{CS02-2}. I was debating whether to learn C++ or just stick with
Python that I am already familiar with; however, I decided to learn
C++ because it never hurts to learn a new language and it is also a
recommended language \cite{CS02-3}. With that in mind, let's get
started with basic syntax from C++!

\section{Foundations}

Let's start with the Hello World example. However to do so, we need
a text editor or an integrated development environment (IDE) and a
compiler. There are many different options for text editor or IDE.
There are online editors that lets you write and compile there.
However, I highly recommend using local editors and compilers for
maintainability. Among many local editors, some popular choices
are Visual Studio Code, CLion, and Vim. Please don't get mad for not
including your favorite editor as it could get very philosophical.
I personally use Neovim due to its low learning curve and ease in
set up. Choosing an editor is up to you!

For compilers, two popular options are GNU Compiler Collection (GCC)
and Clang. I personally use Clang, but this too is up to you!
If you have an editor and a compiler of your choice, let's get
started with printing ``Hello, World!'' on our screen.

\begin{lstlisting}[style=cppstyle]
#include <iostream>

int main() {
    std::cout << "Hello, World!\n";
    return 0;
}
\end{lstlisting}

I know this looks a lot, but let's first compile and see it before
breaking down each line. To compile this, we can run

\begin{lstlisting}[style=basic]
$ clang++ program.cpp -o program
\end{lstlisting}

and execute it with the following commend.

\begin{lstlisting}[style=basic]
$ ./program
\end{lstlisting}

Then, you should see the following output in your terminal.

\begin{lstlisting}[style=basic]
Hello, World!
\end{lstlisting}

Now let's break down the program line by line. First,
\verb|#include <iostream>| is for including the standard input and
output streams library. We need this to print our output on the
screen.

Next, \verb|int main()| defines our \verb|main| function for
execution. Next, \verb|std::cout << "Hello,| \verb|World!\n";| is to
\verb|cout|, or character out, from the standard library specified
with the namespace \verb|std|. We also have \verb|\n| to add a line
break at the end to make a pretty output. Finally, don't forget your
semicolon! We need to tell our compiler that our statement has ended
with the semicolon. Meaning, unlike Python, indentation is only for
writing a legible code.

Lastly, \verb|return 0;| is used to tell that the program has
successfully terminated. Nowadays, it isn't strictly necessary as
reaching the end of the \verb|main| function automatically returns
$0$. However, it's still a good habit to keep them at the end of the
\verb|main| function.

Below is an example with many outputs.

\begin{lstlisting}[style=cppstyle]
#include <iostream>
using namespace std;

int main() {
    cout << "I love banana. " << "Banana likes me" << endl;
    cout << "I ate five bananas today\n";
}
\end{lstlisting}

Compiling and executing returns the following.

\begin{lstlisting}[style=basic]
I love banana. Banana likes me
I ate five bananas today
\end{lstlisting}

Notice that we can add \verb|using namespace std;| to avoid writing
\verb|std::| every time we need \verb|cout|. We can do this for
small projects and speed in competitive programming. However, this
is highly discouraged in header files for large projects as it causes
collisions. One more thing to note is that both \verb|\n| and
\verb|endl| returns a new line. However, they are \textbf{not} the
same thing! We will return to this later when we discuss about
optimization.

\subsection{Variables}

We know how to print a value. How about storing them? Let's first
start with definitions.

\begin{definition}
A region of a storage that a hardware can store a value is called an
\textbf{object}. \textbf{Variables} are objects with a name.
The \textbf{data type} is the type of value that can be stored in the
object \cite{CS02-1}.
\end{definition}

There are many data types and we need to specify the data type
before declaring a variable. Below are some examples of basic
data types.

\begin{center}
\begin{tabularx}{\textwidth}{
    |>{\hsize=.3\hsize\linewidth=\hsize}X
    |>{\hsize=.7\hsize\linewidth=\hsize}X|
}
    \hline
    \textbf{Types} & \textbf{Stores...} \\
    \hline
    int            & Integers \\
    \hline
    long long      & Larger 64-bit integers \\
    \hline
    float          & 4-bit decimals \\
    \hline
    double         & 8-bit decimals \\
    \hline
    char           & Single character \\
    \hline
    string         & Series of characters \\
    \hline
    bool           & True or false value \\
    \hline
\end{tabularx}
\end{center}

With the data type in mind, we can declare a variable as the
following.

\begin{lstlisting}[style=cppstyle]
dataType variableName = value;
\end{lstlisting}

Below is a quick example of declaring a variable and printing
the value.

\begin{lstlisting}[style=cppstyle]
#include <iostream>

int main() {
    std::string name = "Enoch";
    int numBananas = 12 + 34;

    std::cout << "Hi! I am " << name << ", and I ate "
              << numBananas << " bananas today.\n";

    return 0;
}
\end{lstlisting}

Notice that it is important to use quotation when declaring a
string, and we can use simple operations when declaring a variable.
This time, let's take a look at another example where we can
take in the value of a variable and declare variables of the same
type.

\begin{lstlisting}[style=cppstyle]
#include <iostream>

int main() {
    int a, b;

    std::cout << "Type two numbers to add: ";
    std::cin >> a >> b;
    std::cout << "Your sum is: " << a + b << std::endl;
    
    return 0;
}
\end{lstlisting}

When you compile and run it, we can get a simple calculator that
can add two numbers with adequate size (32-bit).

\begin{lstlisting}[style=basic]
Type two numbers to add: 12345 67890
Your sum is: 80235
\end{lstlisting}

If you don't want the value of $a$ to change, we can use \verb|const|.

\begin{lstlisting}[style=cppstyle]
#include <iostream>

int main() {
    const int a = 5;
    int b;

    std::cout << "Type a number: ";
    std::cin >> b;
    std::cout << b << " + 5 = " << a + b << std::endl;
    
    return 0;
}
\end{lstlisting}

In this variation, $a$ is fixed as $5$ and cannot be changed.

\begin{lstlisting}[style=basic]
Type a number: 123456789
123456789 + 5 = 123456794
\end{lstlisting}

Now that we can take inputs and print outputs with various data types,
let's discuss about loops and conditions.

\end{document}


